% GENERATED FILE. Edit canonical lessons, then run npm run generate:books.
% canonical-source-hash: fnv1a64:032591ebbb8d6db0
% canonical-lessons: MW-C262-jano, MW-C262-banjaro, MW-C262-saheli, MW-C262-chaudhri, MW-C262-karigar

\chapter{People of the Village}
\label{ch:mw-people-of-the-village}
\hlchaptermodality{262}

\begin{chapteropening}
\textbf{By the end of this chapter:} \emph{I can name people in the village and say who they are.}

Complete the last lesson of chapter 262: I can name people in the village and say who they are.
\end{chapteropening}

\section[\texorpdfstring{\mw{जणो} (jaṇo)}{जणो (jaṇo)}]{\mw{जणो} (jaṇo) — a person (when counting)}

\label{lesson:MW-C262-jano}



\begin{quote}
Before the new one: say the Marwadi for a son's son, then the Marwadi for a daughter's son.
\end{quote}

\subsection*{You'll want to know: \mw{जणो}}
\textbf{\mw{जणो}} — \emph{jaṇo} — \textquotedblleft{}a person (when counting)\textquotedblright{}.

\begin{quote}
\mw{घर} \mw{मांय} \mw{पांच} \mw{जणा} \mw{है।}
\end{quote}

\begin{quote}
\emph{ghar mā̃y pāṅch jaṇā hai.}
\end{quote}

\begin{quote}
There are five people in the house.
\end{quote}

A masculine noun.

\begin{cousinweb}
\textbf{\mw{जणो}} (\emph{jaṇo}) is inherited from Sanskrit \emph{jana}, \textquotedblleft{}person\textquotedblright{}; CDIAL lists Marwari \emph{jaṇo} under it (CDIAL 5098).
\end{cousinweb}

\begin{grammarlens}[title={What you've built}]
One more word for this chapter.
\end{grammarlens}

\subsection*{Your turn}
\begin{itemize}
\raggedright
  \item \emph{Say it:} \emph{jaṇo}
  \item \emph{Say it:} \emph{jaṇo}, once more
  \item \emph{Say it:} \emph{ghar mā̃y pāṅch jaṇā hai.} — There are five people in the house.
  \item \emph{Recall:} say the Marwadi for a son's son, then the Marwadi for a daughter's son, then say \emph{jaṇo} again
\end{itemize}

\begin{tcolorbox}[breakable,colback=teal!4,colframe=teal!35!black,title={Before you move on}]
What does \mw{जणो} mean? (\textquotedblleft{}A person (when counting)\textquotedblright{}.) Say it once more.
\end{tcolorbox}
\section[\texorpdfstring{\mw{बणजारो} (baṇjāro)}{बणजारो (baṇjāro)}]{\mw{बणजारो} (baṇjāro) — a travelling trader}

\label{lesson:MW-C262-banjaro}



\begin{quote}
Before the new one: say the Marwadi for a daughter's son, then the Marwadi for a person (when counting).
\end{quote}

\subsection*{You'll want to know: \mw{बणजारो}}
\textbf{\mw{बणजारो}} — \emph{baṇjāro} — \textquotedblleft{}a travelling trader\textquotedblright{}.

\begin{quote}
\mw{बणजारो} \mw{ऊंट} \mw{माथै} \mw{लूण} \mw{लावै} \mw{है।}
\end{quote}

\begin{quote}
\emph{baṇjāro ū̃ṭ māthai lūṇ lāvai hai.}
\end{quote}

\begin{quote}
The travelling trader brings salt on a camel.
\end{quote}

A masculine noun.

\begin{cousinweb}
\textbf{\mw{बणजारो}} (\emph{baṇjāro}) is inherited from a reconstructed \emph{*vaṇijyākāra}, \textquotedblleft{}trader\textquotedblright{}. CDIAL gives no Marwari form in this entry; it lists Gujarati \emph{vaṇjāro}, \textquotedblleft{}travelling trader\textquotedblright{}, with the same \mw{ण} (CDIAL 11234).
\end{cousinweb}

\begin{grammarlens}[title={What you've built}]
One more word for this chapter.
\end{grammarlens}

\subsection*{Your turn}
\begin{itemize}
\raggedright
  \item \emph{Say it:} \emph{baṇjāro}
  \item \emph{Say it:} \emph{baṇjāro}, once more
  \item \emph{Say it:} \emph{baṇjāro ū̃ṭ māthai lūṇ lāvai hai.} — The travelling trader brings salt on a camel.
  \item \emph{Recall:} say the Marwadi for a daughter's son, then the Marwadi for a person (when counting), then say \emph{baṇjāro} again
\end{itemize}

\begin{tcolorbox}[breakable,colback=teal!4,colframe=teal!35!black,title={Before you move on}]
What does \mw{बणजारो} mean? (\textquotedblleft{}A travelling trader\textquotedblright{}.) Say it once more.
\end{tcolorbox}
\section[\texorpdfstring{\mw{सहेली} (sahelī)}{सहेली (sahelī)}]{\mw{सहेली} (sahelī) — a woman's friend}

\label{lesson:MW-C262-saheli}



\begin{quote}
Before the new one: say the Marwadi for a person (when counting), then the Marwadi for a travelling trader.
\end{quote}

\subsection*{You'll want to know: \mw{सहेली}}
\textbf{\mw{सहेली}} — \emph{sahelī} — \textquotedblleft{}a woman's friend\textquotedblright{}.

\begin{quote}
\mw{म्हारी} \mw{सहेली} \mw{मेंहदी} \mw{लगावै} \mw{है।}
\end{quote}

\begin{quote}
\emph{mhārī sahelī mẽhdī lagāvai hai.}
\end{quote}

\begin{quote}
My friend is putting on henna.
\end{quote}

A feminine noun.

\begin{cousinweb}
\textbf{\mw{सहेली}} (\emph{sahelī}) is inherited from Sanskrit \emph{sakhī}, \textquotedblleft{}a woman's confidante\textquotedblright{}; CDIAL lists Old Marwari \emph{sahelī} under it (CDIAL 13074).
\end{cousinweb}

\begin{grammarlens}[title={What you've built}]
One more word for this chapter.
\end{grammarlens}

\subsection*{Your turn}
\begin{itemize}
\raggedright
  \item \emph{Say it:} \emph{sahelī}
  \item \emph{Say it:} \emph{sahelī}, once more
  \item \emph{Say it:} \emph{mhārī sahelī mẽhdī lagāvai hai.} — My friend is putting on henna.
  \item \emph{Recall:} say the Marwadi for a person (when counting), then the Marwadi for a travelling trader, then say \emph{sahelī} again
\end{itemize}

\begin{tcolorbox}[breakable,colback=teal!4,colframe=teal!35!black,title={Before you move on}]
What does \mw{सहेली} mean? (\textquotedblleft{}A woman's friend\textquotedblright{}.) Say it once more.
\end{tcolorbox}
\section[\texorpdfstring{\mw{चौधरी} (chaudhrī)}{चौधरी (chaudhrī)}]{\mw{चौधरी} (chaudhrī) — a village headman}

\label{lesson:MW-C262-chaudhri}



\begin{quote}
Before the new one: say the Marwadi for a travelling trader, then the Marwadi for a woman's friend.
\end{quote}

\subsection*{You'll want to know: \mw{चौधरी}}
\textbf{\mw{चौधरी}} — \emph{chaudhrī} — \textquotedblleft{}a village headman\textquotedblright{}.

\begin{quote}
\mw{चौधरी} \mw{घणो} \mw{भलो} \mw{आदमी} \mw{है।}
\end{quote}

\begin{quote}
\emph{chaudhrī ghaṇo bhalo ādmī hai.}
\end{quote}

\begin{quote}
The headman is a very kind man.
\end{quote}

A masculine noun.

\begin{cousinweb}
\textbf{\mw{चौधरी}} (\emph{chaudhrī}) is inherited from Sanskrit \emph{caturdhara}, a family name. CDIAL gives no Marwari form in this entry; it lists Hindi \emph{caudharī}, \textquotedblleft{}the headman of a village\textquotedblright{} (CDIAL 4612).
\end{cousinweb}

\begin{grammarlens}[title={What you've built}]
One more word for this chapter.
\end{grammarlens}

\subsection*{Your turn}
\begin{itemize}
\raggedright
  \item \emph{Say it:} \emph{chaudhrī}
  \item \emph{Say it:} \emph{chaudhrī}, once more
  \item \emph{Say it:} \emph{chaudhrī ghaṇo bhalo ādmī hai.} — The headman is a very kind man.
  \item \emph{Recall:} say the Marwadi for a travelling trader, then the Marwadi for a woman's friend, then say \emph{chaudhrī} again
\end{itemize}

\begin{tcolorbox}[breakable,colback=teal!4,colframe=teal!35!black,title={Before you move on}]
What does \mw{चौधरी} mean? (\textquotedblleft{}A village headman\textquotedblright{}.) Say it once more.
\end{tcolorbox}
\section[\texorpdfstring{\mw{कारीगर} (kārīgar)}{कारीगर (kārīgar)}]{\mw{कारीगर} (kārīgar) — a craftsman; a skilled worker}

\label{lesson:MW-C262-karigar}



\begin{quote}
Before the new one: say the Marwadi for a woman's friend, then the Marwadi for a village headman.
\end{quote}

\subsection*{You'll want to know: \mw{कारीगर}}
\textbf{\mw{कारीगर}} — \emph{kārīgar} — \textquotedblleft{}a craftsman; a skilled worker\textquotedblright{}.

\begin{quote}
\mw{गांव} \mw{मांय} \mw{दो} \mw{कारीगर} \mw{है।}
\end{quote}

\begin{quote}
\emph{gā̃v mā̃y do kārīgar hai.}
\end{quote}

\begin{quote}
There are two craftsmen in the village.
\end{quote}

A masculine noun.

\begin{cousinweb}
\textbf{\mw{कारीगर}} (\emph{kārīgar}) is usually traced to Persian \emph{kārīgar}, \textquotedblleft{}craftsman\textquotedblright{}.
\end{cousinweb}

\begin{grammarlens}[title={What you've built}]
One more word for this chapter.
\end{grammarlens}

\subsection*{Your turn}
\begin{itemize}
\raggedright
  \item \emph{Say it:} \emph{kārīgar}
  \item \emph{Say it:} \emph{kārīgar}, once more
  \item \emph{Say it:} \emph{gā̃v mā̃y do kārīgar hai.} — There are two craftsmen in the village.
  \item \emph{Recall:} say the Marwadi for a woman's friend, then the Marwadi for a village headman, then say \emph{kārīgar} again
\end{itemize}

\begin{tcolorbox}[breakable,colback=teal!4,colframe=teal!35!black,title={Before you move on}]
What does \mw{कारीगर} mean? (\textquotedblleft{}A craftsman; a skilled worker\textquotedblright{}.) Say it once more.
\end{tcolorbox}
